\section{AIME I 1994}
        \begin{enumerate}
        	\item (\textit{AIME I 1994}) The increasing sequence $3, 15, 24, 48, \ldots\,$ consists of those positive multiples of 3 that are one less than a perfect square.  What is the remainder when the 1994th term of the sequence is divided by 1000?


 
 \textbf{Answer: }063



\textbf{Solution}

One less than a perfect square can be represented by $n^2 - 1 = (n+1)(n-1)$. Either $n+1$ or $n-1$ must be divisible by 3. This is true when $n \equiv -1,\ 1 \equiv 2,\ 1 \pmod{3}$. Since 1994 is even, $n$ must be congruent to $1 \pmod{3}$. It will be the $\frac{1994}{2} = 997$th such term, so $n = 4 + (997-1) \cdot 3 = 2992$. The value of $n^2 - 1 = 2992^2 - 1 \pmod{1000}$ is $\boxed{063}$.

 


	\item (\textit{AIME I 1994}) A circle with diameter $\overline{PQ}\,$ of length 10 is internally tangent at $P^{}_{}$ to a circle of radius 20. Square $ABCD\,$ is constructed with $A\,$ and $B\,$ on the larger circle, $\overline{CD}\,$ tangent at $Q\,$ to the smaller circle, and the smaller circle outside $ABCD\,$. The length of $\overline{AB}\,$ can be written in the form $m + \sqrt{n}\,$, where $m\,$ and $n\,$ are integers. Find $m + n\,$.


 
 \textbf{Answer: }312



\textbf{Solution}


                    \begin{figure}[H]
                        \centering
                        \includegraphics[width=0.5\textwidth]{Solution/AIME_I_1994_2_Solution_Figure_1.png}
                    \end{figure}
                    

 Call the center of the larger circle $O$. Extend the diameter $\overline{PQ}$ to the other side of the square (at point $E$), and draw $\overline{AO}$. We now have a \href{https://artofproblemsolving.com/wiki/index.php?title=Right\_triangle}{right triangle}, with \href{https://artofproblemsolving.com/wiki/index.php?title=Hypotenuse}{hypotenuse} of length $20$. Since $OQ = OP - PQ = 20 - 10 = 10$, we know that $OE = AB - OQ = AB - 10$. The other leg, $AE$, is just $\frac 12 AB$. 

 Apply the \href{https://artofproblemsolving.com/wiki/index.php?title=Pythagorean\_Theorem}{Pythagorean Theorem}: 

 
            \begin{enumerate}[label=(\roman*)]
                \item $(AB - 10)^2 + \left(\frac 12 AB\right)^2 = 20^2$\item $AB^2 - 20 AB + 100 + \frac 14 AB^2 - 400 = 0$\item $AB^2 - 16 AB - 240 = 0$
            \end{enumerate}
            The \href{https://artofproblemsolving.com/wiki/index.php?title=Quadratic\_formula}{quadratic formula} shows that the answer is $\frac{16 \pm \sqrt{16^2 + 4 \cdot 240}}{2} = 8 \pm \sqrt{304}$. Discard the negative root, so our answer is $8 + 304 = \boxed{312}$.

 


	\item (\textit{AIME I 1994}) The function $f_{}^{}$ has the property that, for each real number $x,\,$


 $f(x)+f(x-1) = x^2\,$

 .


 If $f(19)=94,\,$ what is the remainder when $f(94)\,$ is divided by 1000?


 
 \textbf{Answer: }561



\textbf{Solution 1}


 \begin{align*}f(94)&=94^2-f(93)=94^2-93^2+f(92)=94^2-93^2+92^2-f(91)=\cdots \\ &= (94^2-93^2) + (92^2-91^2) +\cdots+ (22^2-21^2)+ 20^2-f(19) \\ &= 94+93+\cdots+21+400-94  \\ &= 4561 \end{align*}


 So, the remainder is $\boxed{561}$.

 



\textbf{Solution 2}

Those familiar with triangular numbers and some of their properties will quickly recognize the equation given in the problem. It is well-known (and easy to show) that the sum of two consecutive triangular numbers is a perfect square; that is,
 \[T_{n-1} + T_n = n^2,\]where $T_n = 1+2+...+n = \frac{n(n+1)}{2}$ is the $n$th triangular number.

 Using this, as well as using the fact that the value of $f(x)$ directly determines the value of $f(x+1)$ and $f(x-1),$ we conclude that $f(n) = T_n + K$ for all odd $n$ and $f(n) = T_n - K$ for all even $n,$ where $K$ is a constant real number.

 Since $f(19) = 94$ and $T_{19} = 190,$ we see that $K = -96.$ It follows that $f(94) = T_{94} - (-96) = \frac{94\cdot 95}{2} + 96 = 4561,$ so the answer is $\boxed{561}$.

 


	\item (\textit{AIME I 1994}) Find the positive integer $n\,$ for which


 $\lfloor \log_2{1}\rfloor+\lfloor\log_2{2}\rfloor+\lfloor\log_2{3}\rfloor+\cdots+\lfloor\log_2{n}\rfloor=1994$

 .


 (For real $x\,$, $\lfloor x\rfloor\,$ is the greatest integer $\le x.\,$)


 
 \textbf{Answer: }312



\textbf{Solution 1}

Note that if $2^x \le a<2^{x+1}$ for some $x\in\mathbb{Z}$, then $\lfloor\log_2{a}\rfloor=\log_2{2^{x}}=x$. 

 Thus, there are $2^{x+1}-2^{x}=2^{x}$ integers $a$ such that $\lfloor\log_2{a}\rfloor=x$. So the sum of $\lfloor\log_2{a}\rfloor$ for all such $a$ is $x\cdot2^x$. 

 Let $k$ be the integer such that $2^k \le n<2^{k+1}$. So for each integer $j<k$, there are $2^j$ integers $a\le n$ such that $\lfloor\log_2{a}\rfloor=j$, and there are $n-2^k+1$ such integers such that $\lfloor\log_2{a}\rfloor=k$. 

 Therefore, $\lfloor\log_2{1}\rfloor+\lfloor\log_2{2}\rfloor+\lfloor\log_2{3}\rfloor+\cdots+\lfloor\log_2{n}\rfloor= \sum_{j=0}^{k-1}(j\cdot2^j) + k(n-2^k+1) = 1994$. 

 Through computation: $\sum_{j=0}^{7}(j\cdot2^j)=1538<1994$ and $\sum_{j=0}^{8}(j\cdot2^j)=3586>1994$. Thus, $k=8$. 

 So, $\sum_{j=0}^{k-1}(j\cdot2^j) + k(n-2^k+1) = 1538+8(n-2^8+1)=1994 \Rightarrow n = \boxed{312}$.

 Alternatively, one could notice this is an \href{https://artofproblemsolving.com/wiki/index.php?title=Arithmetico-geometric\_series}{arithmetico-geometric series} and avoid a lot of computation.

 



\textbf{Solution 2}

For this solution, we notice that values between ranges of $2^n$ are the same. Let's look at these ranges and their total value. It is trivial to conclude $\log_2{1} = 0$ so we will not write that.

 \begin{array}{c|c}\text{Range of } k & \text{Expression (Total)} \\\hline2 \leq k < 2^2 & 1 \cdot (2^2 - 2) = 2 \\2^2 \leq k < 2^3 & 2 \cdot (2^3 - 2^2) = 8 \\2^3 \leq k < 2^4 & 3 \cdot (2^4 - 2^3) = 24 \\2^4 \leq k < 2^5 & 4 \cdot (2^5 - 2^4) = 64 \\2^5 \leq k < 2^6 & 5 \cdot (2^6 - 2^5) = 160 \\2^6 \leq k < 2^7 & 6 \cdot (2^7 - 2^6) = 384 \\2^7 \leq k < 2^8 & 7 \cdot (2^8 - 2^7) = 896 \\\end{array}

 Sum of the values up to $1$ less than $2^k$ is $1538$. This seems close enough to our desired value of $1994$ as any additional groups would exceed $1994$. The difference is $1996 - 1538 = 456$. Since the next numbers of the sequence will always be 8 since $\lfloor \log_2{2^8 + n < 2^9} \rfloor = 8$, we can just find the number of '8's, which is $\frac{456}{8} = 57$. So there are exactly 57 integers in that range. The largest integer, which is $a$, is equal to $2^8 + 56 = \boxed{312}$. ~Totient4Breakfast

 



\textbf{Solution 3 (Variation of Solution 1)}

Like in solution 1, continue upon realization that we get the series \( 0 + 2(1) + 4(2) + 8(3) + \dots + k(n) \).

 Then, this series can be written as
 \[\sum_{k=0}^n k2^k = 1994\]Continuing forth with solution 1, we can see that if \( n=7 \), we get \( 1538 < 1994 \), and if \( n=8 \), we get \( 3586 > 1994 \).

 Then, we see that there must be \( m(8) \) such that \( 1538 + 8m = 1994 \). Solving for \( m \) we get \( m=57 \).

 We see that \( 2^k + 2^{k+1} + 2^{k+2} + \dots + 2^{k+n} \) is the number of terms in this sequence for \( 0 \le k \le 7 \) and \( n=7 \), and we include the additional 57 terms. To do this, we don't want to include the arithmetic portion of our series, turning our arithmetico-geometric series into
 \[\left(\sum_{k=0}^n 2^k\right) + 57\]The finite geometric sum results in 255, in which the value of \( n \) is determined as \( 255 + 57 = \)  $\boxed{312}$.

 


	\item (\textit{AIME I 1994}) Given a positive integer $n\,$, let $p(n)\,$ be the product of the non-zero digits of $n\,$.  (If $n\,$ has only one digit, then $p(n)\,$ is equal to that digit.)  Let


 $S=p(1)+p(2)+p(3)+\cdots+p(999)$

 .


 What is the largest prime factor of $S\,$?


 
 \textbf{Answer: }103



\textbf{Solution 1}

Suppose we write each number in the form of a three-digit number (so $5 \equiv 005$), and since our $p(n)$ ignores all of the zero-digits, replace all of the $0$s with $1$s. Now note that in the expansion of 

 $(1+ 1 +2+3+4+5+6+7+8+9) (1+ 1 +2+3+\cdots+9) (1+ 1 +2+3+\cdots+9)$

 we cover every permutation of every product of $3$ digits, including the case where that first $1$ represents the replaced $0$s. However, since our list does not include $000$, we have to subtract $1$. Thus, our answer $\boxed{103}$.

 



\textbf{Solution 2}

Note that $p(1)=p(11), p(2)=p(12), p(3)=p(13), \cdots p(19)=p(9)$, and $p(37)=3p(7)$. So $p(10)+p(11)+p(12)+\cdots +p(19)=46$, $p(10)+p(11)+\cdots +p(99)=46*45=2070$. We add $p(1)+p(2)+p(3)+\cdots +p(10)=45$ to get 2115. When we add a digit we multiply the sum by that digit. Thus $2115\cdot (1+1+2+3+4+5+6+7+8+9)=2115\cdot 46=47\cdot 45\cdot 46$. But we didn't count 100, 200, 300, ..., 900. Thus, we have to add another 45 to get $45\cdot 2163$. The largest prime factor of that is $\boxed{103}$.

 


	\item (\textit{AIME I 1994}) The graphs of the equations


 $y=k, \qquad y=\sqrt{3}x+2k, \qquad y=-\sqrt{3}x+2k,$

 are drawn in the coordinate plane for $k=-10,-9,-8,\ldots,9,10.\,$  These 63 lines cut part of the plane into equilateral triangles of side $2/\sqrt{3}$.  How many such triangles are formed?


 
 \textbf{Answer: }660



\textbf{Solution 1}

We note that the lines partition the hexagon of the six extremal lines into disjoint unit regular triangles, and forms a series of unit regular triangles along the edge of the hexagon. 

 
                    \begin{figure}[H]
                        \centering
                        \includegraphics[width=0.5\textwidth]{Solution/AIME_I_1994_6_Solution_Figure_1.png}
                    \end{figure}
                    

 Solving the above equations for $k=\pm 10$, we see that the hexagon in question is regular, with side length $\frac{20}{\sqrt{3}}$. Then, the number of triangles within the hexagon is simply the ratio of the area of the hexagon to the area of a regular triangle. Since the ratio of the area of two similar figures is the square of the ratio of their side lengths, we see that the ratio of the area of one of the six equilateral triangles composing the regular hexagon to the area of a unit regular triangle is just $\left(\frac{20/\sqrt{3}}{2/\sqrt{3}}\right)^2 = 100$. Thus, the total number of unit triangles is $6 \times 100 = 600$.

 There are $6 \cdot 10$ equilateral triangles formed by lines on the edges of the hexagon. Thus, our answer is $600+60 = \boxed{660}$.

 



\textbf{Solution 2}

There are three types of lines: horizontal, upward-slanting diagonal, and downward-slanting diagonal. There are $21$ of each type of line, and a unit equilateral triangle is determined by exactly one of each type of line. Given an upward-slanting diagonal and a downward-slanting diagonal, they determine exactly two unit equilateral triangles as shown below.
                    \begin{figure}[H]
                        \centering
                        \includegraphics[width=0.5\textwidth]{Solution/AIME_I_1994_6_Solution_Figure_2.png}
                    \end{figure}
                    Therefore, if all horizontal lines are drawn, there will be a total of $2\cdot 21^2=882$ unit equilateral triangles. Of course, we only draw $21$ horizontal lines, so we are overcounting the triangles that are caused by the undrawn horizontal lines. In the below diagram, we draw the diagonal lines and the highest and lowest horizontal lines.
                    \begin{figure}[H]
                        \centering
                        \includegraphics[width=0.5\textwidth]{Solution/AIME_I_1994_6_Solution_Figure_3.png}
                    \end{figure}
                    

 We see that the lines $y=-21,-20,\dots, -11$ and $y=11,12,\dots,21$ would complete several of the $882$ unit equilateral triangles. In fact, we can see that the lines $y=-21,-20,\dots,-11$ complete $1,2,(1+3),(2+4),(3+5),(4+6),\dots,(9+11)$ triangles, or $111$ triangles. The positive horizontal lines complete the same number of triangles, hence the answer is $882-2\cdot 111=\boxed{660}$.

 


 



\textbf{Solution 3 Elementary Counting}

Picturing the diagram in your head should give you an illustration similar to the one above. The distance from parallel sides of the center hexagon is 20 units, and by extending horizontal lines to the sides of the hexagon. This shows that for every side of the hexagon there are 10 spaces. Therefore the side length of the biggest triangle (imagine one of the overlapping triangles in the Star of David) is 30. The total number of triangles in the hexagon can be found by finding the number of triangles in the extended triangle and subtracting the 3 corner triangles. This gives us $30^2 - 10^2 - 10^2- 10^2 = 600$. That is the number of triangles in the hexagon. The remaining triangles form in groups of 10 on the exterior of each side. $600 + 6 * 10 = \boxed{660}$.

 
-jackshi2006

 


	\item (\textit{AIME I 1994}) For certain ordered pairs $(a,b)\,$ of real numbers, the system of equations


 $ax+by=1\,$

 $x^2+y^2=50\,$

 has at least one solution, and each solution is an ordered pair $(x,y)\,$ of integers.  How many such ordered pairs $(a,b)\,$ are there?


 
 \textbf{Answer: }072



\textbf{Solution}

The equation $x^2+y^2=50$ is that of a circle of radius $\sqrt{50}$, centered at the origin. By testing integers until the left side becomes too big, we see that the lattice points on this circle are $(\pm1,\pm7)$, $(\pm5,\pm5)$, and $(\pm7,\pm1)$ where the signs are all independent of each other, for a total of $3\cdot 2\cdot 2=12$ lattice points. They are indicated by the blue dots below.

 
                    \begin{figure}[H]
                        \centering
                        \includegraphics[width=0.5\textwidth]{Solution/AIME_I_1994_7_Solution_Figure_1.png}
                    \end{figure}
                    

 Since $(x,y)=(0,0)$ yields $a\cdot 0+b\cdot 0=0 \neq 1$, we know that $ax+by=1$ is the equation of a line that does not pass through the origin. So, we are looking for the number of lines which pass through at least one of the $12$ lattice points on the circle, but do not pass through the origin or through any non-lattice point on the circle. An example is the green line above. It is straightforward to show that a line passes through the origin precisely when there exist two opposite points $(p,q)$ and $(-p,-q)$ through which it passes. And example is the red line above.

 There are $\binom{12}{2}=66$ ways to pick two distinct lattice points, and subsequently $66$ distinct lines which pass through two distinct lattice points on the circle. Then we subtract the lines which pass through the origin by noting that the lattice points on the circle can be grouped into opposite pairs $(p,q)$ and $(-p,-q)$, for a total of $\frac{12}{2}=6$ lines. Finally, we add the $12$ unique tangent lines to the circle at each of the lattice points.

 Therefore, our final count of distinct lines which pass through one or two of the lattice points on the circle, but do not pass through the origin, is
 \[66-6+12=\boxed{72}.\]

 


	\item (\textit{AIME I 1994}) The points $(0,0)\,$, $(a,11)\,$, and $(b,37)\,$ are the vertices of an equilateral triangle.  Find the value of $ab\,$.


 
 \textbf{Answer: }315



\textbf{Solution 1}

Consider the points on the \href{https://artofproblemsolving.com/wiki/index.php?title=Complex\_plane}{complex plane}. The point $b+37i$ is then a rotation of $60$ degrees of $a+11i$ about the origin, so:

 
 \[(a+11i)\left(\mathrm{cis}\,60^{\circ}\right) = (a+11i)\left(\frac 12+\frac{\sqrt{3}i}2\right)=b+37i.\]

 Equating the real and imaginary parts, we have:

 
 \begin{align*}b&=\frac{a}{2}-\frac{11\sqrt{3}}{2}\\37&=\frac{11}{2}+\frac{a\sqrt{3}}{2} \end{align*}


 Solving this system, we find that $a=21\sqrt{3}, b=5\sqrt{3}$. Thus, the answer is $\boxed{315}$.

 \textbf{Note}: There is another solution where the point $b+37i$ is a rotation of $-60$ degrees of $a+11i$. However, this triangle is just a reflection of the first triangle by the $y$-axis, and the signs of $a$ and $b$ are flipped. The product $ab$ is unchanged.

 



\textbf{Solution 2}

Using the Pythagorean theorem with these beastly numbers doesn't seem promising. How about properties of equilateral triangles? $\sqrt{3}$ and perpendiculars inspires this solution:

 First, drop a perpendicular from $O$ to $AB$. Call this midpoint of $AB$ $M$. Thus, $M=\left(\frac{a+b}{2}, 24\right)$. The vector from $O$ to $M$ is $\left[\frac{a+b}{2}, 24\right]$. Meanwhile from point $M$ we can use a vector with $\frac{\sqrt{3}}{3}$ the distance; we have to switch the $x$ and $y$ and our displacement is $\left[8\sqrt{3}, \frac{(a+b)\sqrt{3}}{6}\right]$. (Do you see why we switched $x$ and $y$ due to the rotation of 90 degrees?)

 
We see this displacement from $M$ to $A$ is $\left[\frac{a-b}{2}, 13\right]$ as well. Equating the two vectors, we get $a+b=26\sqrt{3}$ and $a-b=16\sqrt{3}$. Therefore, $a=21\sqrt{3}$ and $b=5\sqrt{3}$. And the answer is $\boxed{315}$.

 \textbf{Note}: This solution was also present in Titu Andreescu and Zuming Feng's "103 Trigonometry Problems".

 



\textbf{Solution 3}

Plot this equilateral triangle on the complex plane.Translate the equilateral triangle so that its centroid is located at the origin. (The centroid can be found by taking the average of the three vertices of the triangle, which gives $\left(\frac{a+b}{3}, 16i\right)$. The new coordinates of the equilateral triangle are $\left(-\frac{a+b}{3}-16i\right), \left(a-\frac{a+b}{3}-5i\right), \left(b-\frac{a+b}{3}+21i\right)$. These three vertices are solutions of a cubic polynomial of form $x^3 + C$. By Vieta's Formulas, the sum of the paired roots of the cubic polynomial are zero. (Or for the three roots $r_1, r_2,$ and $r_3,$$\, r_1r_2 + r_2r_3 + r_3r_1 = 0$.) The vertices of the equilateral triangle represent the roots of a polynomial, so the vertices can be plugged into the above equation. Because both the real and complex components of the equation have to sum to zero, you really have two equations. Multiply out the equation given by Vieta's Formulas and isolate the ones with imaginary components. Simplify that equation, and that gives the equation $5a = 21b.$Now use the equation with only real parts. This should give you a quadratic $a^2 - ab + b^2 = 1083$. Use your previously obtained equation to plug in for $a$ and solve for $b$, which should yield $5\sqrt{3}$. $a$ is then $\frac{21}{5}\sqrt{3}$. Multiplying $a$ and $b$ yields $\boxed{315}$.

 



\textbf{Solution 4 (gigantic numbers)}

Just using the Pythagorean Theorem, we get that $a^2 + 11^2 = (b-a)^2 + 26^2 = b^2 + 37^2$. 

 $a^2 + 121 = b^2 + 1369 \longrightarrow a^2 = b^2 + 1248$. Expanding the second and subtracting the first equation from it we get $b^2 = 2ab - 555$.

 $b^2 = 2ab - 555 \longrightarrow a^2 = 2ab + 693$. 

 We have $b^2 + 1248 = 2b\sqrt{b^2+1248} + 693$. 

 Moving the square root to one side and non square roots to the other we eventually get $b^4 + 1110b^2 + 308025 = 4b^4 + 4992b^2$. 

 $3b^4 + 3882b^2 - 308025 = 0$. 

 This factors to $(3b^2-225)(b^2+1369)$, so $3b^2 = 225, b = 5\sqrt 3$. 

 Plugging it back in, we find that a = $\sqrt{1323}$ which is $21\sqrt3$, so the product $ab$ is $\boxed {315}$.

 


	\item (\textit{AIME I 1994}) A solitaire game is played as follows.  Six distinct pairs of matched tiles are placed in a bag.  The player randomly draws tiles one at a time from the bag and retains them, except that matching tiles are put aside as soon as they appear in the player's hand.  The game ends if the player ever holds three tiles, no two of which match; otherwise the drawing continues until the bag is empty.  The probability that the bag will be emptied is $p/q,\,$ where $p\,$ and $q\,$ are relatively prime positive integers.  Find $p+q.\,$


 
 \textbf{Answer: }394



\textbf{Solution}

Let $P_k$ be the \href{https://artofproblemsolving.com/wiki/index.php?title=Probability}{probability} of emptying the bag when it has $k$ pairs in it. Let's consider the possible draws for the first three cards:

 
                \begin{itemize}
                    \item Case 1. We draw a pair on the first two cards. The second card is the same as the first with probability $\frac {1}{2k - 1}$, then we have $k - 1$ pairs left. So this contributes probability $\frac {P_{k - 1}}{2k - 1}$.\item Case 2. We draw a pair on the first and third cards. The second card is different from the first with probability $\frac {2k - 2}{2k - 1}$ and the third is the same as the first with probability $\frac {1}{2k - 2}$. We are left with $k - 1$ pairs but one card already drawn. However, having drawn one card doesn't affect the game, so this also contributes probability $\frac {P_{k - 1}}{2k - 1}$.\item Case 3. We draw a pair on the second and third cards. This is pretty much the same as case 2, so we get $\frac {P_{k - 1}}{2k - 1}$.
                \end{itemize}


                Therefore, we obtain the \href{https://artofproblemsolving.com/wiki/index.php?title=Recursion}{recursion} $P_k = \frac {3}{2k - 1}P_{k - 1}$. Iterating this for $k = 6,5,4,3,2$ (obviously $P_1 = 1$), we get $\frac {3^5}{11*9*7*5*3} = \frac {9}{385}$, and $p+q=\boxed{394}$.

 



\textbf{Solution 2}

Call the case that we begin with [ABCDEF]. It doesn't matter what letter we choose at first, so WLOG assume we choose A. Now there is BCDEFABCDEF remaining in the bag. We have two cases to consider here.

 1. We pick the other A. There's a $\frac{1}{11}$ chance for this to happen. We remain with the case [BCDEF] if this is the case.

 2. We pick any other letter that is not an A. There's a $\frac{10}{11}$ chance for this to happen. WLOG, assume we pick the letter B. Now in order for us to continue the game, we must choose either the other A or B. There's a $\frac{2}{10}$ chance for this to happen. WLOG, assume we choose A. Now we have BCDEFCDEF left.

 Notice however that in the first case, the probability of emptying the bag with [BCDEF] is the same thing as with BCDEFCDEF, as the only difference is you've removed one of the letters (and it doesn't matter which you chose). 

 Hence for this case, there is a $\frac{1}{11} + \frac{10}{11}*\frac{2}{10} = \frac{3}{11}$ * [BCDEF] chance to empty the bag.

 Continuing this process, we get that: 

 [BCDEF] = $\frac{3}{9}$ * [CDEF]

 [CDEF] = $\frac{3}{7}$ * [DEF]

 [DEF] = $\frac{3}{5}$ * [EF]

 [EF] = 1 (clearly, since if we are only left with EFEF then we are going to empty the bag).

 And hence [ABCDEF] = $1*\frac{3}{5}*\frac{3}{7}*\frac{3}{9}*\frac{3}{11} = \frac{9}{385}$ so our answer is $9+385=\boxed{394}$

 



\textbf{Solution 3}

Let $a_{m,n}$ represent the situation in which $n$ tiles are in your hand and $m$ pairs have been removed; i.e., there are $2m$ tiles remaining in play. We know that $a_{m,0}=a_{m,1}$ to begin with, so we find some other recursions:

 When the player has $1$ tile in their hand with $m$ pairs remaining, there is a $\frac{1}{2m-1}$ chance that the player finds the one matching tile and the game state changes to $a_{m-1,0}$. There is a $\frac{2m-2}{2m-1}$ chance that the player does not find the matching tile and the game state changes to $a_{m,2}$. As a result, $a_{m,1}=\frac{1}{2m-1}a_{m-1,0}+\frac{2m-2}{2m-1}a_{m,2}$.

 When the player has $2$ tiles in their hand with $m$ pairs remaining, there is a $\frac{1}{m-1}$ chance that the player finds a matching tile and the game moves on to game state $a_{m-1,1}$. As a result, $a_{m,2}=\frac{1}{m-1}a_{m-1,1}$.

 We also know that $a_0=a_1=a_2=1$ (i.e., you are guaranteed to win with two tiles remaining), so we are able to find an end to the recursion.

 Now we are left with four rules:

 1) $a_{m,0}=a_{m,1}$

 2) $a_{m,1}=\frac{1}{2m-1}a_{m-1,0}+\frac{2m-2}{2m-1}a_{m,2}$

 3) $a_{m,2}=\frac{1}{m-1}a_{m-1,1}$

 4) $a_0=a_1=a_2=1$

 We then calculate the terms of the multi-variable recursion and arrive at $a_{6,0}=\frac{9}{385}$, so the answer is $385+9=\boxed{394}$.

 


	\item (\textit{AIME I 1994}) In triangle $ABC,\,$ angle $C$ is a right angle and the altitude from $C\,$ meets $\overline{AB}\,$ at $D.\,$  The lengths of the sides of $\triangle ABC\,$ are integers, $BD=29^3,\,$ and $\cos B=m/n\,$, where $m\,$ and $n\,$ are relatively prime positive integers.  Find $m+n.\,$


 
 \textbf{Answer: }450



\textbf{Solution 1}

Since $\triangle ABC \sim \triangle CBD$, we have $\frac{BC}{AB} = \frac{29^3}{BC} \Longrightarrow BC^2 = 29^3 AB$. It follows that $29^2 | BC$ and $29 | AB$, so $BC$ and $AB$ are in the form $29^2 x$ and $29 x^2$, respectively, where $x$ is an integer. 

 By the \href{https://artofproblemsolving.com/wiki/index.php?title=Pythagorean\_Theorem}{Pythagorean Theorem}, we find that $AC^2 + BC^2 = AB^2 \Longrightarrow (29^2x)^2 + AC^2 = (29 x^2)^2$, so $29x | AC$. Letting $y = AC / 29x$, we obtain after dividing through by $(29x)^2$, $29^2 = x^2 - y^2 = (x-y)(x+y)$. As $x,y \in \mathbb{Z}$, the pairs of factors of $29^2$ are $(1,29^2)(29,29)$; clearly $y = \frac{AC}{29x} \neq 0$, so $x-y = 1, x+y= 29^2$. Then, $x = \frac{1+29^2}{2} = 421$.

 Thus, $\cos B = \frac{BC}{AB} = \frac{29^2 x}{29x^2} = \frac{29}{421}$, and $m+n = \boxed{450}$.

 



\textbf{Solution 2}

We will solve for $\cos B$ using $\triangle CBD$, which gives us $\cos B = \frac{29^3}{BC}$. By the Pythagorean Theorem on $\triangle CBD$, we have $BC^2 - DC^2 = (BC + DC)(BC - DC) = 29^6$. Trying out factors of $29^6$, we can either guess and check or just guess to find that $BC + DC = 29^4$ and $BC - DC = 29^2$ (The other pairs give answers over 999). Adding these, we have $2BC = 29^4 + 29^2$ and $\frac{29^3}{BC} = \frac{2*29^3}{29^2 (29^2 +1)} = \frac{58}{842} = \frac{29}{421}$, and our answer is $\boxed{450}$.

 



\textbf{Solution 3}

Using similar right triangles, we identify that $CD = \sqrt{AD \cdot BD}$. Let $AD$ be $29 \cdot k^2$, to avoid too many radicals, getting $CD = k \cdot 29^2$. Next we know that $AC = \sqrt{AB \cdot AD}$ and that $BC = \sqrt{AB \cdot BD}$. Applying the logic with the established values of k, we get $AC = 29k \cdot \sqrt{29^2 + k^2}$ and $BC = 29^2 \cdot \sqrt{29^2 + k^2}$. Next we look to the integer requirement. Since $k$ is both outside and inside square roots, we know it must be an integer to keep all sides as integers. Let $y$ be $\sqrt{29^2 + k^2}$, thus $29^2 = y^2 - k^2$, and $29^2 = (y + k)(y - k)$. Since $29$ is prime, and $k$ cannot be zero, we find $k = 420$ and $y = 421$ as the smallest integers satisfying this quadratic Diophantine equation. Then, since $cos B$ = $\frac{29}{\sqrt{29^2 + k^2}}$. Plugging in we get $cos B = \frac{29}{421}$, thus our answer is $\boxed{450}$.

 


	\item (\textit{AIME I 1994}) Ninety-four bricks, each measuring $4''\times10''\times19'',$ are to be stacked one on top of another to form a tower 94 bricks tall.  Each brick can be oriented so it contributes $4''\,$ or $10''\,$ or $19''\,$ to the total height of the tower.  How many different tower heights can be achieved using all 94 of the bricks?


 
 \textbf{Answer: }465



\textbf{Solution 1}

We have the smallest stack, which has a height of $94 \times 4$ inches. Now when we change the height of one of the bricks, we either add $0$ inches, $6$ inches, or $15$ inches to the height. Now all we need to do is to find the different change values we can get from $94$ $0$'s, $6$'s, and $15$'s. Because $0$, $6$, and $15$ are all multiples of $3$, the change will always be a multiple of $3$, so we just need to find the number of changes we can get from $0$'s, $2$'s, and $5$'s.

 From here, we count what we can get:

 
 \[0, 2 = 2, 4 = 2+2, 5 = 5, 6 = 2+2+2, 7 = 5+2, 8 = 2+2+2+2, 9 = 5+2+2, \ldots\]

 It seems we can get every integer greater or equal to four; we can easily deduce this by considering \href{https://artofproblemsolving.com/wiki/index.php?title=Parity}{parity} or using the \href{https://artofproblemsolving.com/wiki/index.php?title=Chicken\_McNugget\_Theorem}{Chicken McNugget Theorem}, which says that the greatest number that cannot be expressed in the form of $2m + 5n$ for $m,n$ being \href{https://artofproblemsolving.com/wiki/index.php?title=Positive\_integer}{positive integers} is $5 \times 2 - 5 - 2=3$. 

 But we also have a maximum change ($94 \times 5$), so that will have to stop somewhere. To find the gaps, we can work backwards as well. From the maximum change, we can subtract either $0$'s, $3$'s, or $5$'s. The maximum we can't get is $5 \times 3-5-3=7$, so the numbers $94 \times 5-8$ and below, except $3$ and $1$, work. Now there might be ones that we haven't counted yet, so we check all numbers between $94 \times 5-8$ and $94 \times 5$. $94 \times 5-7$ obviously doesn't work, $94 \times 5-6$ does since 6 is a multiple of 3, $94 \times 5-5$ does because it is a multiple of $5$ (and $3$), $94 \times 5-4$ doesn't since $4$ is not divisible by $5$ or $3$, $94 \times 5-3$ does since $3=3$, and $94 \times 5-2$ and $94 \times 5-1$ don't, and $94 \times 5$ does.

 Thus the numbers $0$, $2$, $4$ all the way to $94 \times 5-8$, $94 \times 5-6$, $94 \times 5-5$, $94 \times 5-3$, and $94\times 5$ work. That's $2+(94 \times 5 - 8 - 4 +1)+4=\boxed{465}$ numbers. That's the number of changes you can make to a stack of bricks with dimensions $4 \times 10 \times 19$, including not changing it at all.

 



\textbf{Solution 2}

Using bricks of dimensions $4''\times10''\times19''$ is comparable to using bricks of dimensions $0''\times6''\times15''$ which is comparable to using  bricks of dimensions $0''\times2''\times5''$. Using 5 bricks of height $2''$ can be replaced by using 2 bricks of height $5''$ and 3 bricks of height $0''$. 

 It follows that all tower heights can be made by using 4 or fewer bricks of height $2''$. There are $95+94+93+92+91=465$ ways to build a tower using 4 or fewer bricks of height $2''$. Taking the heights $\mod 5$, we see that towers using a different number of bricks of height $2''$ have unequal heights. Thus, all of the $\boxed{465}$ tower heights are different.

 


	\item (\textit{AIME I 1994}) A fenced, rectangular field measures 24 meters by 52 meters. An agricultural researcher has 1994 meters of fence that can be used for internal fencing to partition the field into congruent, square test plots. The entire field must be partitioned, and the sides of the squares must be parallel to the edges of the field. What is the largest number of square test plots into which the field can be partitioned using all or some of the 1994 meters of fence? 


 
 \textbf{Answer: }702



\textbf{Solution 1}

Suppose there are $n$ squares in every column of the grid, so there are $\frac{52}{24}n = \frac {13}6n$ squares in every row. Then $6|n$, and our goal is to maximize the value of $n$. 

 Each vertical fence has length $24$, and there are $\frac{13}{6}n - 1$ vertical fences; each horizontal fence has length $52$, and there are $n-1$ such fences. Then the total length of the internal fencing is $24\left(\frac{13n}{6}-1\right) + 52(n-1) = 104n - 76 \le 1994 \Longrightarrow n \le \frac{1035}{52} \approx 19.9$, so $n \le 19$. The largest multiple of $6$ that is $\le 19$ is $n = 18$, which we can easily verify works, and the answer is $\frac{13}{6}n^2 = \boxed{702}$.

 



\textbf{Solution 2}

Assume each partitioned square has a side length of $1$ (just so we can get a clear image of what the formula will look like). The amount of internal fencing that is required to partition the field is clearly $52*(24+1) + 24(52+1)$. (If you are confused, just draw the square out). This is clearly greater than $1994$, so the actual side length that we are looking for is greater than $1$.

 Now we can convert this into an equation. The equation is simply $(\frac{52}{x})(\frac{24}{x}+1) + (\frac{24}{x})(\frac{52}{x}+1)$ (The intutition comes from considering partioning the field into side lengths of $1$ and then partitioning those squares). This is equivalent to $\frac{2496}{x^2} + \frac{76}{x}$, which should be less than or equal to $1994$.

 Now we can just find possible lengths of the square that are greater than $1$ and test them out. A viable side length would mean that $\frac{24}{x}, \frac{52}{x} \in$ N. Since $\gcd(24,52) = 4$, then the smallest value greater than $1$ that we satisfies the conditions has $4$ in the numerator, and hence $3$ in the denominator. Test out $x=\frac{4}{3}$. This will equate to something less than $1994$, and hence the smallest square length that is plausible is $\frac{4}{3}$. 

 Now the rest is elementary, we do $\frac{52}{\frac{4}{3}} * \frac{24}{\frac{4}{3}} \Rightarrow 39*18 = \boxed{702}$

 


	\item (\textit{AIME I 1994}) The equation


 $x^{10}+(13x-1)^{10}=0\,$

 has 10 complex roots $r_1, \overline{r_1}, r_2, \overline{r_2}, r_3, \overline{r_3}, r_4, \overline{r_4}, r_5, \overline{r_5},\,$ where the bar denotes complex conjugation.  Find the value of


 $\frac 1{r_1\overline{r_1}}+\frac 1{r_2\overline{r_2}}+\frac 1{r_3\overline{r_3}}+\frac 1{r_4\overline{r_4}}+\frac 1{r_5\overline{r_5}}.$

 
 \textbf{Answer: }850



\textbf{Solution 1}

Let $t = 1/x$. After multiplying the equation by $t^{10}$, $1 + (13 - t)^{10} = 0\Rightarrow (13 - t)^{10} = - 1$.

 Using DeMoivre, $13 - t = \text{cis}\left(\frac {(2k + 1)\pi}{10}\right)$ where $k$ is an integer between $0$ and $9$.

 $t = 13 - \text{cis}\left(\frac {(2k + 1)\pi}{10}\right) \Rightarrow \bar{t} = 13 - \text{cis}\left(-\frac {(2k + 1)\pi}{10}\right)$.

 Since $\text{cis}(\theta) + \text{cis}(-\theta) = 2\cos(\theta)$, $t\bar{t} = 170 - 26\cos \left(\frac {(2k + 1)\pi}{10}\right)$ after expanding. Here $k$ ranges from 0 to 4 because two angles which sum to $2\pi$ are involved in the product.

 The expression to find is $\sum t\bar{t} = 850 - 26\sum_{k = 0}^4 \cos \frac {(2k + 1)\pi}{10}$.

 But $\cos \frac {\pi}{10} + \cos \frac {9\pi}{10} = \cos \frac {3\pi}{10} + \cos \frac {7\pi}{10} = \cos \frac {\pi}{2} = 0$ so the sum is $\boxed{850}$.

 



\textbf{Solution 2}

Divide both sides by $x^{10}$ to get 
 \[1 + \left(13-\dfrac{1}{x}\right)^{10}=0\]

 Rearranging: 
 \[\left(13-\dfrac{1}{x}\right)^{10} = -1\]

 Thus, $13-\dfrac{1}{x} = \omega$ where $\omega = e^{i(\pi n/5+\pi/10)}$ where $n$ is an integer.

 We see that $\dfrac{1}{x}=13-\omega$. Thus, 
 \[\dfrac{1}{x\overline{x}}=(13\, -\, \omega)(13\, -\, \overline{\omega})=169-13(\omega\, +\, \overline{\omega})\, +\, \omega\overline{\omega}=170\, -\, 13(\omega\, +\, \overline{\omega})\]

 Summing over all terms: 
 \[\dfrac{1}{r_1\overline{r_1}}+\cdots + \dfrac{1}{r_5\overline{r_5}} = 5\cdot 170 - 13(e^{i\pi/10}+\cdots +e^{i(9\pi/5+\pi/10)})\]

 However, note that $e^{i\pi/10}+\cdots +e^{i(9\pi/5+\pi/10)}=0$ from drawing the numbers on the complex plane, our answer is just 
 \[5\cdot 170=\boxed{850}\]

 



\textbf{Solution 3(Ileytyn)}

Let us apply difference of squares, by writing this equation as $(x^5)^2-(i[13x-1]^5)^2=0$, where $i=\sqrt{-1}$. 

 Once applied, we have 
 \[(x^5-i[13x-1]^5)(x^5+i[13x-1]^5)=0\] By factorization, we get that either 
 \[\text{equation 1:   }x^5-i[13x-1]^5=0\text{ or equation 2:    }x^5+i[13x-1]^5=0\] We find the trivial solution to the first equation $x^5=i[13x-1]^5$, and since a fifth root of $i=i^5$ is $i$, we can find this solution by taking the fifth root, or when $x=(13x-1)i$, which when solved, gives $x=\frac{i}{13i-1}=\frac{13-i}{170}$. By using a similar approach to the second equation, $x^5+i[13x-1]^5=0$, or $x=-i(13x-1)$, we get $\frac{i}{13i+1}=\frac{13+i}{170}$. We observe that these are conjugates. We also note that $z\overline{z}=|z|^2$. Hence, we find the magnitude of one of these two conjugates, which is $\frac{1}{\sqrt{170}}$, or as we are finding the reciprocal of the square value of this, this contributes $170$ to the final answer. We claim that each of the roots of equation one is a conjugate of one of equation 2. We also notice that, by an application of DeMoivre's, each of the roots has the same magnitude, just a different angle. Hence, as each of the reciprocals of product of conjugates has the value $170$, and there are 5 of these pairs of conjugates, the answer is 
 \[170\cdot 5 = \boxed{850}\]

 


	\item (\textit{AIME I 1994}) A beam of light strikes $\overline{BC}\,$ at point $C\,$ with angle of incidence $\alpha=19.94^\circ\,$ and reflects with an equal angle of reflection as shown.  The light beam continues its path, reflecting off line segments $\overline{AB}\,$ and $\overline{BC}\,$ according to the rule: angle of incidence equals angle of reflection.  Given that $\beta=\alpha/10=1.994^\circ\,$ and $AB=BC,\,$ determine the number of times the light beam will bounce off the two line segments.  Include the first reflection at $C\,$ in your count.


 
                    \begin{figure}[H]
                        \centering
                        \includegraphics[width=0.5\textwidth]{Problem/AIME_I_1994_14_Problem_Figure_1.png}
                    \end{figure}
                    


 
 \textbf{Answer: }071



\textbf{Solution}

At each point of reflection, we pretend instead that the light continues to travel straight. 
                    \begin{figure}[H]
                        \centering
                        \includegraphics[width=0.5\textwidth]{Solution/AIME_I_1994_14_Solution_Figure_1.png}
                    \end{figure}
                    Note that after $k$ reflections (excluding the first one at $C$) the extended line will form an angle $k \beta$ at point $B$. For the $k$th reflection to be just inside or at point $B$, we must have $k\beta \le 180 - 2\alpha \Longrightarrow k \le \frac{180 - 2\alpha}{\beta} = 70.27$. Thus, our answer is, including the first intersection, $\left\lfloor \frac{180 - 2\alpha}{\beta} \right\rfloor + 1 = \boxed{071}$.

 


	\item (\textit{AIME I 1994}) Given a point $P^{}_{}$ on a triangular piece of paper $ABC,\,$ consider the creases that are formed in the paper when $A, B,\,$ and $C\,$ are folded onto $P.\,$  Let us call $P_{}^{}$ a fold point of $\triangle ABC\,$ if these creases, which number three unless $P^{}_{}$ is one of the vertices, do not intersect.  Suppose that $AB=36, AC=72,\,$ and $\angle B=90^\circ.\,$  Then the area of the set of all fold points of $\triangle ABC\,$ can be written in the form $q\pi-r\sqrt{s},\,$ where $q, r,\,$ and $s\,$ are positive integers and $s\,$ is not divisible by the square of any prime.  What is $q+r+s\,$?


 
 \textbf{Answer: }597



\textbf{Solution}

Let $O_{AB}$ be the intersection of the \href{https://artofproblemsolving.com/wiki/index.php?title=Perpendicular\_bisector}{perpendicular bisectors} (in other words, the intersections of the creases) of $\overline{PA}$ and $\overline{PB}$, and so forth. Then $O_{AB}, O_{BC}, O_{CA}$ are, respectively, the \href{https://artofproblemsolving.com/wiki/index.php?title=Circumcenter}{circumcenters} of $\triangle PAB, PBC, PCA$. According to the problem statement, the circumcenters of the triangles cannot lie within the interior of the respective triangles, since they are not on the paper. It follows that $\angle APB, \angle BPC, \angle CPA > 90^{\circ}$; the \href{https://artofproblemsolving.com/wiki/index.php?title=Locus}{locus} of each of the respective conditions for $P$ is the region inside the (semi)circles with diameters $\overline{AB}, \overline{BC}, \overline{CA}$.

 We note that the circle with diameter $AC$ covers the entire triangle because it is the circumcircle of $\triangle ABC$, so it suffices to take the intersection of the circles about $AB, BC$. We note that their intersection lies entirely within $\triangle ABC$ (the chord connecting the endpoints of the region is in fact the altitude of $\triangle ABC$ from $B$). Thus, the area of the locus of $P$ (shaded region below) is simply the sum of two \href{https://artofproblemsolving.com/wiki/index.php?title=Segment}{segments} of the circles. If we construct the midpoints of $M_1, M_2 = \overline{AB}, \overline{BC}$ and note that $\triangle M_1BM_2 \sim \triangle ABC$, we see that thse segments respectively cut a $120^{\circ}$ arc in the circle with radius $18$ and $60^{\circ}$ arc in the circle with radius $18\sqrt{3}$. 

 
                    \begin{figure}[H]
                        \centering
                        \includegraphics[width=0.5\textwidth]{Solution/AIME_I_1994_15_Solution_Figure_1.png}
                    \end{figure}
                    

 The diagram shows $P$ outside of the grayed locus; notice that the creases [the dotted blue] intersect within the triangle, which is against the problem conditions. The area of the locus is the sum of two segments of two circles; these segments cut out $120^{\circ}, 60^{\circ}$ angles by simple similarity relations and angle-chasing.

 Hence, the answer is, using the $\frac 12 ab\sin C$ definition of triangle area, $\left[\frac{\pi}{3} \cdot 18^2 - \frac{1}{2} \cdot 18^2 \sin \frac{2\pi}{3} \right] + \left[\frac{\pi}{6} \cdot \left(18\sqrt{3}\right)^2 - \frac{1}{2} \cdot (18\sqrt{3})^2 \sin \frac{\pi}{3}\right] = 270\pi - 324\sqrt{3}$, and $q+r+s = \boxed{597}$.

 


\end{enumerate}